\begin{corollary}[5.11.4]
\EGAReaderTarget{ega4.c04.cor.5-11-4}{IV.2.s5.11.4-fr}
Sean $X$ un preesquema localmente noetheriano, $Z$ una parte cerrada de $X$,
$U=X\setmin Z$, $i:U\to X$ la inyección canónica, $\mathcal F$ un
$\mathcal O_U$-módulo coherente. Para que $i_*(\mathcal F)$ sea un
$\mathcal O_X$-módulo coherente, es necesario que para todo
$x\in\operatorname{Ass}(\mathcal F)$ se tenga
$\operatorname{codim}(\overline{\{x\}}\cap Z,\overline{\{x\}})\geq2$.
Esta condición es suficiente en cada uno de los dos casos siguientes:
\begin{enumerate}
\item[\textup{(i)}] El preesquema $X$ es localmente inmersible en un esquema regular.
\item[\textup{(ii)}] El preesquema $X$ es universalmente catenario y universalmente
japonés \textup{(lo que significa, por definición, que todo punto $x\in X$ admite una
vecindad abierta afín cuyo anillo es} universalmente japonés \textup{(\textbf{0},
23.1.1))}.
\end{enumerate}
\end{corollary}

Se sabe (\textbf I, 9.4.7) que existe un sub-$\mathcal O_X$-módulo \emph{coherente}
$\mathcal G$ de $i_*(\mathcal F)$ tal que $\mathcal G\vert U=\mathcal F$. Se
tiene evidentemente
$\operatorname{Ass}(\mathcal F)\subset\operatorname{Ass}(\mathcal G)
\subset\operatorname{Ass}(i_*(\mathcal F))$, y como
\[
\operatorname{Ass}(i_*(\mathcal F))=\operatorname{Ass}(\mathcal F)
\]
(3.1.13), se tiene $\operatorname{Ass}(\mathcal G)=\operatorname{Ass}(\mathcal F)$;
basta entonces aplicar (5.11.1) al $\mathcal O_X$-módulo coherente
$\mathcal G$, observando que $i_*(\mathcal F)=\mathcal H^0_{X/Z}(\mathcal G)$
y que, cuando $X$ es universalmente catenario y universalmente japonés,
la hipótesis (ii) de (5.11.1) se verifica por definición.

\begin{corollary}[5.11.5]
\label{IV.2.s5.11.5-fr}
Sean $X$ un preesquema localmente noetheriano, $Z$ una parte de $X$ estable por
especialización, $\mathcal F$ un $\mathcal O_X$-módulo coherente tal que se
tenga $\operatorname{Ass}(\mathcal F)\subset X\setmin Z$. Entonces la condición:
\begin{enumerate}
\item[\textup{a)}] $\mathcal H^0_{X/Z}(\mathcal F)$ es un $\mathcal O_X$-módulo
coherente
\end{enumerate}
implica la siguiente:
\begin{enumerate}
\item[\textup{d)}] Para toda parte $T$ de $Z$ cerrada en $X$ \textup{(o solamente para una}
familia filtrante

\oldpage[IV-2]{125}

creciente \textup{de partes cerradas $T$ de $Z$, de reunión $Z$)},
$i_*(\mathcal F\vert X\setmin T)$ (donde $i:X\setmin T\to X$ es la inyección canónica) es
coherente.
\end{enumerate}
Cuando $X$ verifica una de las hipótesis \textup{(i), (ii)} de \textup{(5.11.1), a)} y \textup{d)} son
equivalentes.
\end{corollary}

Observemos que se tiene
$\operatorname{Ass}(i_*(\mathcal F\vert X\setmin T))=\operatorname{Ass}(\mathcal F)$
en virtud de la hipótesis y de (3.1.13); se sigue, pues, de (5.10.2) que las
aplicaciones canónicas
\[
i_*(\mathcal F\vert X\setmin T)\longrightarrow\mathcal H^0_{X/Z}(\mathcal F)
\]
son \emph{inyectivas}; el hecho de que \emph{a)} implique \emph{d)} es, por tanto, consecuencia de
esta observación. Recíprocamente, la condición \emph{d)} implica, en virtud de
(5.11.1), que $\operatorname{codim}(T\cap Y_\alpha,Y_\alpha)\geq2$ con las
notaciones de (5.11.1); por consiguiente, se tiene
$\operatorname{codim}(Z\cap Y_\alpha,Y_\alpha)\geq2$ puesto que $Z$ es reunión
de sus partes que son cerradas en $X$, y la última afirmación del
corolario se deduce de (5.11.1).

\begin{corollary}[5.11.6]
\label{IV.2.s5.11.6-fr}
Sean $A$ un anillo noetheriano, $X=\operatorname{Spec}(A)$. Consideremos las
propiedades siguientes:
\begin{enumerate}
\item[\textup{a)}] Para todo anillo cociente íntegro $B$ de $A$, el anillo $B^{(1)}$
\textup{(notación de (5.11.2))} es una $B$-álgebra finita.
\item[\textup{b)}] Para todo $\mathcal O_X$-módulo coherente $\mathcal F$ y toda parte
$Z$ de $X$, estable por especialización, y tal que para todo
$x\in\operatorname{Ass}(\mathcal F)\cap(X\setmin Z)$ se tenga
$\operatorname{codim}(\overline{\{x\}}\cap Z,\overline{\{x\}})\geq2$, el
$\mathcal O_X$-módulo $\mathcal H^0_{X/Z}(\mathcal F)$ es coherente.
\item[\textup{c)}] Para toda parte cerrada $T$ de $X$ y todo $\mathcal O_U$-módulo
coherente $\mathcal G$ (donde $U=X\setmin T$) tal que para todo
$x\in\operatorname{Ass}(\mathcal G)$ se tenga
$\operatorname{codim}(\overline{\{x\}}\cap T,\overline{\{x\}})\geq2$,
$i_*(\mathcal G)$ (donde $i:U\to X$ es la inyección canónica) es un
$\mathcal O_X$-módulo coherente.
\item[\textup{d)}] Para todo anillo cociente íntegro $B$ de $A$ y todo ideal
$\mathfrak J$ de altura $\geq2$ en $B$, el anillo
$\bigcap_{\mathfrak p\not\supset\mathfrak J}B_\mathfrak p$
es una $B$-álgebra finita.
\end{enumerate}
Se tienen entonces las implicaciones
\[
\mathrm a)\Longleftrightarrow\mathrm b)\Longrightarrow\mathrm c)
\Longleftrightarrow\mathrm d)
\]
Además, las condiciones \textup{a), b), c)} y \textup{d)} se verifican en cada uno de los dos
casos siguientes:
\begin{enumerate}
\item[\textup{(i)}] $A$ es cociente de un anillo regular.
\item[\textup{(ii)}] $A$ es universalmente catenario y universalmente japonés.
\end{enumerate}
\end{corollary}

Sea $B=A/\mathfrak q$, donde $\mathfrak q$ es un ideal primo de $A$, de modo
que $Y=\operatorname{Spec}(B)$ es la parte cerrada $V(\mathfrak q)$ de $X$;
pongamos $Z=Z^{(2)}(Y)$, que es una parte de $X$ estable por especialización;
como $B$ es íntegro, $\operatorname{Ass}(A/\mathfrak q)$ se reduce al punto
genérico $\mathfrak q$ de $Y$. Si se verifica la condición \emph{b)}, se puede
aplicarla al $\mathcal O_X$-módulo coherente
$\mathcal F=(A/\mathfrak q)^\sim$ y a $Z$, y, en virtud de (5.9.3.1), esto
muestra que $B^{(1)}$ es un $A$-módulo de tipo finito, y \emph{a fortiori} un
$B$-módulo de tipo finito. Recíprocamente, supongamos verificada \emph{a)}; entonces, si
$\mathcal F$ es un $\mathcal O_X$-módulo coherente tal que para todo
$x\in\operatorname{Ass}(\mathcal F)\cap(X\setmin Z)$ se tenga
$\operatorname{codim}(\overline{\{x\}}\cap Z,\overline{\{x\}})\geq2$, se puede
aplicar (con las notaciones de (5.11.1)) el resultado de \emph{a)} a cada uno de los esquemas afines
$Y_\alpha=\operatorname{Spec}(B_\alpha)$, donde $B_\alpha$ es un anillo cociente
íntegro de $A$; como por hipótesis $Z_\alpha$ está contenido
en $Z^{(2)}(Y_\alpha)$, la condición \emph{a)} (teniendo en cuenta (5.10.2)

\oldpage[IV-2]{126}

y el hecho de que $\operatorname{Ass}(\mathcal O_{Y_\alpha})$ se reduce al punto
genérico de $Y_\alpha$) implica que
$\mathcal H^0_{Y_\alpha/Z_\alpha}(\mathcal O_{Y_\alpha})$ es un
$\mathcal O_{Y_\alpha}$-módulo coherente, y \emph{b)} se deduce entonces de (5.11.1).

Para ver que \emph{c)} implica \emph{d)}, se razona como anteriormente aplicando \emph{c)} al
caso en que $\mathcal F=(A/\mathfrak q)^\sim$ y $Z=V(\mathfrak J)$;
recíprocamente, se prueba que \emph{d)} implica \emph{c)} utilizando de nuevo la equivalencia de
\emph{a)} y \emph{b)} en (5.11.1). Es evidente que \emph{c)} es un caso particular de \emph{b)}.
Finalmente, el hecho de que \emph{a)} (y, por consiguiente, cada una de las demás condiciones) se verifique
cuando $A$ verifica una de las hipótesis (i), (ii) se deduce de (5.11.2), dada la
definición de los anillos universalmente catenarios y de los anillos
universalmente japoneses.

\begin{remarks}[5.11.7]
\label{IV.2.s5.11.7-fr}
\begin{enumerate}
\item[(i)] Se ignora si, en (5.11.5), la condición \emph{d)} implica \emph{a)} sin una
hipótesis adicional sobre $X$; veremos más adelante (7.2.4) que, en efecto, es
así cuando $X$ es un esquema \emph{local}. De igual modo, demostraremos que
cuando $A$ es un anillo \emph{local} noetheriano las cuatro propiedades \emph{a)}, \emph{b)}, \emph{c)},
\emph{d)} de (5.11.6) son equivalentes (7.2.4). Se ignora si este resultado se extiende a
todos los anillos noetherianos.
\item[(ii)] Si $A$ verifica la propiedad \emph{a)} de (5.11.6), lo mismo ocurre con
todo \emph{anillo de fracciones} $S^{-1}A$ y toda $A$-álgebra \emph{finita} $C$. En efecto,
todo anillo cociente íntegro de $S^{-1}A$ es de la forma
$S^{-1}(A/\mathfrak q)$, donde $\mathfrak q$ es un ideal primo de $A$ que no
interseca $S$; y, por otra parte, si $\mathfrak r$ es un ideal primo de
$C$, $\mathfrak p$ su imagen inversa en $A$, $C/\mathfrak r$ es una
$(A/\mathfrak p)$-álgebra finita íntegra que contiene $A/\mathfrak p$; nuestras
afirmaciones son, pues, consecuencias de (5.10.17, (iii) y (iv)).
\end{enumerate}
\end{remarks}
